\section{Introduction}\label{sec:introduction}

We study the Cauchy problem for the hard-sphere Boltzmann equation on
\(\Omega=\T^3_x\times\R^3_v\), where
\(\T^3=\R^3/\mathbb Z^3\). A nonnegative density \(F(t,x,v)\)
evolves by free transport and binary elastic collisions:
\[
 (\partial_t+v\cdot\nabla_x)F=Q(F,F),\qquad F(0,x,v)=F_0(x,v).
\]
For incoming velocities \(v,v_*\) and a collision normal
\(\omega\in S^2\), put
\[
 a=(v-v_*)\cdot\omega,\qquad
 v'=v-a\omega,\qquad v_*'=v_*+a\omega.
\]
Our normalization is
\[
 Q(F,F)(v)=\int_{\R^3\times S^2}|a|
       \bigl(F(v')F(v_*')-F(v)F(v_*)\bigr)\dd v_*\dd\omega,
\]
with the common variables \(t,x\) suppressed and ordinary surface area
measure on \(S^2\). The initial density is assumed to have finite
mass, kinetic energy, and absolute entropy, meaning that
\(\int_\Omega F_0(1+|v|^2+|\log F_0|)<\infty\), with
\(0|\log0|=0\). The question is whether these physical bounds and the
renormalized equation determine the subsequent evolution uniquely.

The global existence and weak stability theory of DiPerna and Lions
provides solutions at this level of integrability
\cite{DiPernaLions1989}. Their entropy-production inequality explicitly
includes periodic space \cite{DiPernaLions1991};
\cite[Theorem~4.1]{LevermoreMasmoudi2010} gives a periodic formulation
with a Maxwellian reference density. The need for renormalization is
substantive: the stated bounds do not themselves make the quadratic
collision term locally integrable. Existence and weak stability do not
supply uniqueness. The uniqueness problem for general renormalized
solutions is discussed, for example, in
\cite[Section~2]{Silvestre2022}.

We use the following precise admissibility requirements, written in
full in Definition~\ref{def:admissible}. With
\(M(v)=(2\pi)^{-3/2}e^{-|v|^2/2}\), the equation is tested through
\(M\beta(F/M)\) for every \(C^1\) function \(\beta\) satisfying
\(\beta(0)=0\) and \(|\beta'(z)|\le C_\beta/(1+z)\); the separately
renormalized gain and loss are integrable on every finite time interval.
The solution is weakly continuous in \(L^1(\Omega)\), against every
bounded measurable test, and takes the prescribed initial value.
Its mass, energy, and absolute entropy remain bounded on bounded time
intervals. It conserves local mass and total momentum and satisfies
both the global energy inequality and the relative-entropy inequality
with the full Boltzmann entropy dissipation at every time.
Here admissibility does not require local momentum or local energy
conservation. Our main result concerns uniqueness within this class.

\begin{theorem}\label{thm:main}
There is a nonnegative density \(F_0\) on \(\T^3\times\R^3\), with bounded
velocity support and
\[
\int_{\T^3\times\R^3}
 F_0(1+|v|^2+|\log F_0|)\dd x\dd v<\infty,
\]
for which the hard-sphere Boltzmann equation has two distinct global
admissible solutions, in the sense of Definition~\ref{def:admissible},
with initial datum \(F_0\).
On some common initial interval both solutions are strongly continuous
in \(L^1\), and their collision gains and losses are integrable in
space-time-velocity.
\end{theorem}

The two solutions therefore obey the unrenormalized equation on the
interval where their difference is established. Their common datum has
bounded velocities but unbounded amplitudes concentrated at successively
smaller spatial scales. The conservation and dissipation requirements
remain part of the conclusion.

\subsection{Uniqueness and related constructions}

Classical uniqueness theorems impose additional control, such as
Maxwellian upper bounds \cite{KanielShinbrot1978}. For periodic cutoff
hard potentials, Ukai's perturbative theory gives global existence and
uniqueness near a Maxwellian in a velocity-weighted spatial Sobolev class,
with mass, momentum, and energy matched to that Maxwellian
\cite[Theorem~1.2]{Ukai1974}.

The distinction between spatially homogeneous and inhomogeneous evolution
is also important. For the homogeneous hard-sphere equation, Mischler and
Wennberg prove uniqueness in a conservative class from finite mass and
energy alone \cite[Theorem~1.1]{MischlerWennberg1999}. Wennberg's
homogeneous nonuniqueness examples for hard potentials permit increasing
energy \cite{Wennberg1999}. Theorem~\ref{thm:main} instead concerns a
spatially inhomogeneous datum, with the energy and entropy-dissipation
inequalities retained.

Recent instability constructions give another relevant comparison.
Chen and Holmer prove failure of uniform continuity in weighted Sobolev
spaces for the constant collision kernel
\cite[Theorem~1.3]{ChenHolmer2024}; Chen, Guo, Shen, and Zhang establish failure of
uniform continuity for the hard-sphere solution map on
\(\R^3_x\times\R^3_v\), in Gaussian-weighted
\(L^2_vH^s_x\) spaces for \(0\le s<1\)
\cite[Theorem~1.1]{ChenGuoShenZhang2026}. Those results compare solutions
with arbitrarily close initial data. Here the two initial data agree
exactly, on periodic space, and both resulting solutions satisfy the
admissibility conditions above. The instability in
\cite{ChenGuoShenZhang2026} is driven by collision loss and dispersion.
In the construction below, collision gain amplifies a vanishing hot seed
and produces a nonzero weak limit near the
initial spatial singularity.

For the related Krieger--Strain model, Gismondi, Golding, and Novack
construct rough stationary solutions in a mollifier-confluent sense:
the collision operators of mollified densities tend to zero in
distributions for every nonnegative Schwartz mollifier of mass one
\cite[Definition~1.3 and Theorem~1.5]{GismondiGoldingNovack2026}.
They also conjecture
nonuniqueness of renormalized inhomogeneous Landau and Boltzmann solutions
for a dense class of initial data. Theorem~\ref{thm:main} concerns one
periodic hard-sphere datum in the specified admissibility class.

\subsection{One datum and two approximation families}

Write \(x=(y,z)\in\R^2\times\R\) in a lift of the torus, and
\(v=(\bar v,v_3)\in\R^2\times\R\). The datum consists of narrow
velocity distributions, called cold jets, near the axial directions
\(\pm e_3\). They occupy transverse annuli of geometrically decreasing
radii \(s_j\downarrow0\). Within each annulus the two velocity signs
occupy disjoint, finely alternating stripes. Their thermal widths are
much smaller than the stripe gaps. An expanding axial velocity profile
bounds the spatial density of each freely transported jet at scale
\(s\) by \(C/(t+s)\). Annular mass is \(O(s^2)\), and annular absolute entropy
is \(O(s^2(1+|\log s|))\), so all scales together define a single
finite-entropy datum.

Assign a color to each cold jet. A collision within one color retains
that color; collisions between different colors, or with particles
already declared hot, produce hot particles. This is an exact
decomposition: the sum of the colored equations is the Boltzmann
equation. The narrow velocity distributions suppress self-collisions,
and the spatial gaps suppress collisions between distinct cold jets
at early times.

Retain the first \(N\) annuli and impose a collision cutoff that is
removed as \(N\to\infty\). Let \(h_N\) denote the hot component.
The estimates use an integrable velocity weight \(H(v)\), Gaussian at
infinity and locally comparable to \(|\bar v|^{-q}\) near the velocity axis,
with \(1<q<2\). A \emph{cap} is an almost-everywhere upper bound on
\(h_N/H\). The two approximation families have the following properties:
\begin{enumerate}
\item With no hot particles initially, the hot density remains bounded
by \(t^2H(v)\). We call this the dormant family.
\item For every sufficiently short observation horizon \(T>0\), a
smooth hot seed of vanishing mass, energy, and entropy can be selected
so that \(h_N\le(t+s_N)H\) on \([0,T]\), with equality in the
weighted essential supremum at some time. Here \(s_N\) is the smallest
retained annular radius.
\end{enumerate}
At each fixed horizon, uniform integrability of both the densities and
their two-particle products gives admissible limits with exactly the
same initial datum. A vanishing seed and a cap hit alone do not imply
that these limits differ: the hit might carry arbitrarily little mass.
Excluding that loss of information is the central step.

\subsection{From a cap hit to separation}

Along a sequence of observation horizons shrinking to zero, rescale
time and space at a cap time \(t_{*,N}\) by \(r_N=t_{*,N}+s_N\),
and divide the hot density by \(r_N\).
The resulting cap is of order one. The cold density has a different
normalization: multiplying it by \(r_N\) produces a locally bounded velocity measure
at positive times. Its two axial velocity signs coexist in the limit because
the increasingly fine spatial stripes average to their occupied
fractions. Strong convergence of hot velocity averages identifies the
cold--hot product. This use of velocity averaging follows the
uniformly integrable formulation of
\cite[Proposition~3]{GolseLionsPerthameSentis1988}.
The limit obeys a linear Boltzmann equation in that two-sign background.
This dilation is distinct from the fixed-horizon compactness used to
construct admissible solutions.

To show that the dilated hot limit is nonzero, we analyze normalized
Carleman gain kernels. Carleman's representation expresses a collision
gain as an integral over velocity planes; its normalization here is
derived in Section~\ref{ker:section}. Iterated-gain regularization has
an earlier role in the homogeneous theory
\cite[Lemma~2.1]{MischlerWennberg1999}. In the present construction the
estimates must be uniform as the cold velocity distributions collapse
onto the axis, despite the singularity of \(H\).
Two complementary arguments provide this control. Near the velocity
axis, two successive gains have vanishing normalized mass: the first
has a uniformly non-atomic axial input, while a second large gain would
force that input to match a small collection of characteristic cold
speeds. Away from the axis, five gains suffice. Their times and angular
parameters map to final time, position, and velocity with rank seven
almost everywhere. A compact submersion argument then turns vanishing
local hot mass into a vanishing gain contribution. Together these facts
prevent the cap hit from disappearing in the weak limit.

The limiting hot equation also has a branching interpretation. A hot
particle colliding with a cold jet creates two hot outputs. We retain
finite collision histories with bounded velocities and show that their
mass grows by a fixed factor inside an expanding spatial box. If the
dilated observation horizon were infinite, this growth would exceed
the polynomial mass bound supplied by the cap. The horizon is therefore
finite, so the cap scale remains comparable to the physical observation
horizon along the extracted sequence.

Finally, the dormant solution vanishes off the velocity axis under the
corresponding shrinking-horizon dilation. If cap limits agreed with it
at arbitrarily small fixed horizons, a diagonal choice of approximants
would force the nonzero dilated hot limit to vanish as well. This
contradiction yields two distinct early solutions of the original
equation with the same datum. The periodic existence theorem then
continues each solution globally while preserving admissibility.
Figure~\ref{fig:proof-route} separates the two compactness procedures
and their roles in this final comparison.

\begin{figure}[p]
\centering
\begin{tikzpicture}[
 box/.style={draw,rounded corners=2pt,align=center,text width=4.8cm,
             inner sep=6pt,font=\small},
 every path/.style={line width=.55pt}, >=Stealth]
\node[box] (datum) at (0,0)
 {One cold-jet datum $F_0$\\finite mass, energy,\\and entropy};
\node[box] (dormant) at (-3.05,-1.8)
 {No hot seed\\$h_N\le t^2H$};
\node[box] (seed) at (3.05,-1.8)
 {Vanishing hot seed\\$h_N\le(t+s_N)H$, with a cap hit};
\node[box] (dl) at (-3.05,-3.7)
 {Fixed-horizon limit\\admissible solution $F^{\mathrm d}$};
\node[box] (sl) at (3.05,-3.7)
 {Fixed-horizon limit\\admissible cap solution};
\node[box] (zero) at (-3.05,-5.9)
 {Shrinking-horizon dilation\\vanishes off the velocity axis};
\node[box] (dilation) at (3.05,-5.9)
 {Dilation of cap approximants\\nonzero hot limit $g$};
\node[box] (growth) at (3.05,-8.0)
 {Branching in the\\limiting background\\finite dilated horizon $U$};
\node[box] (separate) at (0,-10.2)
 {Same-datum solutions differ\\then continue each globally};
\draw[->] (datum) -- (dormant);
\draw[->] (datum) -- (seed);
\draw[->] (dormant) -- (dl);
\draw[->] (seed) -- (sl);
\draw[->] (seed.east) -- (6.05,-1.8) |- (dilation.east);
\draw[->] (dilation) -- (growth);
\draw[->] (dl) -- (zero);
\draw[->] (zero.south) -- (-3.05,-10.2) -- (separate.west);
\draw[->] (sl.west) -- (0,-3.7) -- (separate.north);
\draw[->] (growth.south) -- (3.05,-10.2) -- (separate.east);
\end{tikzpicture}

\caption{Proof structure. At each fixed physical horizon, the two
approximation families yield admissible solutions with the same initial
datum. Separately, cap approximants on shrinking horizons are dilated
about a cap hit. Nonconcentration preserves a nonzero hot limit, and
branching prevents its scale from vanishing relative to the horizon along
the extracted sequence. These facts contradict equality with the dormant
limit, whose corresponding dilation vanishes off the velocity axis.}
\label{fig:proof-route}
\end{figure}

\subsection{Organization}

Section~\ref{sec:framework} gives the full admissibility definition and
the colored equations. Section~\ref{ker:section} establishes uniform
collision-kernel estimates, and Section~\ref{jet:section} constructs the
datum and controls its cold components. Section~\ref{sec:cap} selects
the approximation families; Section~\ref{sec:admissibility} constructs
their fixed-horizon limits and verifies admissibility.
Section~\ref{sec:dilation} identifies the dilated linear equation, and
Section~\ref{nc:section} proves that its hot limit is nonzero.
Section~\ref{sec:branching} excludes an infinite dilated horizon.
Section~\ref{sec:separation} completes the same-datum comparison and
global continuation.
